\documentclass[12pt,a4paper]{article}
\usepackage[utf8]{inputenc}
\usepackage[spanish]{babel}
\usepackage{geometry}
\geometry{top=2.5cm, bottom=2.5cm, left=3cm, right=3cm}
\usepackage{graphicx}
\usepackage{tabularx}
\usepackage{booktabs}
\usepackage{hyperref}
\usepackage{cite}

\title{TA-TAR-01: Aplicación de la Norma IEEE 29148 \\ Sistema Analizado: App Móvil Tuenti Ecuador}
\author{Luis Alberto Paca Pilataza \\ Ingeniería de Requerimientos (ISR-401)}
\date{Noviembre 2026}

\begin{document}

\maketitle

\section{Resumen Analítico: Estándar ISO/IEC/IEEE 29148:2018}

La especificación de requisitos es la etapa fundamental en la ingeniería de software donde se define si el sistema resolverá el problema correcto. De acuerdo con el estándar internacional ISO/IEC/IEEE 29148:2018, un requisito no es simplemente un deseo ambiguo del cliente, sino una declaración documentada que traduce una necesidad en una condición o capacidad medible que el sistema debe poseer.

\subsection{Características de un Requisito Individual Bien Formado}
Según la cláusula 5.2.4 de la norma IEEE 29148, para que un enunciado sea considerado un requerimiento de ingeniería válido, debe cumplir estrictamente con nueve características individuales:

\begin{itemize}
    \item \textbf{Necesario:} La condición especificada debe ser indispensable. Sin ella, el sistema no sería aceptado por los stakeholders.
    \item \textbf{Apropiado al nivel:} El detalle de la redacción debe corresponder al nivel de abstracción del documento (negocio, sistema o software).
    \item \textbf{No ambiguo:} Debe admitir una única interpretación posible por parte de todos los actores involucrados.
    \item \textbf{Completo:} No debe dejar información implícita; debe especificar entradas, comportamiento y salidas.
    \item \textbf{Singular:} Debe contener una sola capacidad o restricción por enunciado para facilitar su trazabilidad.
    \item \textbf{Factible:} Debe poder realizarse considerando los límites tecnológicos, económicos y legales.
    \item \textbf{Verificable:} Si un requisito no se puede comprobar mediante una prueba, inspección o demostración, es técnicamente inaceptable.
    \item \textbf{Correcto:} Refleja con fidelidad la necesidad real expresada por la fuente.
    \item \textbf{Conforme:} Debe respetar el formato, estilo y uso de verbos modales definidos en las plantillas del proyecto.
\end{itemize}



\end{document}
